\uuid{qqjx}
\titre{Noyau/Image/Bijection}
\auteur{hanani ; mzari}
\organisation{exo7}
\licence{CC BY-NC-SA 4.0}
\datecreate{2020-01-20}
\niveau{L1}
\module{Algèbre}
\chapitre{Application linéaire}
\sousChapitre{Définition}
\difficulte{3}
\theme{Noyau et image, Noyau, Image, Bijection}
\exo7qcmid{lille-2.applications-lineaires.033}

\contenu{
\begin{question}
\qclassification{107.01}
\qtags{motcle=Noyau/Image/Bijection}
Soit $E$ un espace vectoriel  et $f$ un endomorphisme de $E$. Quelles sont les assertions vraies ?
\begin{answers}
\bad{Si $f^2 =0$, alors $f=0$}.
\bad{Si $f^2 =0$, alors $f$ est bijective}.
\good{Si $f^2 =0$, alors $\Im f \subset \ker f$}.
\good{Si $\Im f \subset \ker f$, alors  $f^2 =0$}.
\end{answers}
\begin{explanations} Si $f^2 =0$, $f$ n'est pas forcément nulle. Exemple : $f:\Rr^2 \to \Rr^2$ définie par : $f(x,y)=(x-y,x-y)$.
\vskip0mm
Supposons que $f^2 =0$ et que $f\neq 0$. Alors, il existe $x_0\in E$ tel que $f(x_0) \neq 0$ et $f(x_0) \in \ker f$, puisque $f^2(x_0) =0$. Donc $f$ est non injective et donc non bijective.
\vskip0mm
Soit $y\in \Im f$, alors il existe $x\in E$ tel que $y=f(x)$. De l'égalité $f^2 =0$, on déduit que $f(y)=0$.
\vskip0mm
Soit $x\in E$, alors $f(x)\in \Im f$. Comme  $\Im f \subset \ker f$, $f^2(x)=0$.
\end{explanations}
\end{question}
}
